% 晶圆级集成概述
\section{晶圆级集成与晶圆级交换机}

晶圆级集成的核心使能技术是硅interposer——一块无源硅基板，通过微凸点（$\mu$bump）将多个裸片键合在同一平面上。
与PCB互联相比，硅interposer的$\mu$bump pitch可达36--55\,$\mu$m，布线密度高出两个数量级~\cite{ucie2.0-2024}；
每bit传输功耗在UCIe Advanced标准下低至0.25\,pJ，而传统SerDes跨背板方案在10--20\,pJ量级~\cite{oif-cei-5.1}。
TSMC的InFO-SoW和CoWoS平台已使该技术进入量产——
Tesla Dojo通过InFO-SoW集成了25个D1裸片~\cite{tesla2022dojo}，
Cerebras WSE-3以单晶圆级芯片集成约90万个处理核心~\cite{cerebras2022wse2}。

晶圆级平台特别适合构建网络交换机，原因有三。
第一，裸片间互联密度远超传统封装——单个裸片可布置约7万个$\mu$bump，
聚合IO带宽可达Pbps量级，为超大radix交换网络提供了物理基础。
第二，互联协议灵活可配——UCIe Advanced（短距低功耗）、SerDes VSR/MR/LR（中远距）、
光学互联（超远距高带宽）可在同一晶圆上混合使用，根据裸片间距自适应选择~\cite{ucie2.0-2024,oif-cei-5.1}。
第三，交换拓扑可直接映射到interposer上的裸片布局——每个裸片承载若干交换端口，
裸片间物理链路对应拓扑中的边，消除了传统多机箱交换机跨机箱互联的带宽瓶颈。

Chen等人~\cite{chen2024waferscale}首次量化了晶圆级交换机的radix提升潜力；
Feng和Ma~\cite{feng2024switchless}提出利用晶圆级高密度互联替代传统高radix交换机的无交换机Dragonfly架构；
Wan等人~\cite{wan2024architectural}在2D Mesh上叠加逻辑拓扑实现了896端口晶圆级交换系统；
Yang等人~\cite{yang2025pd}提出的TickTock框架面向LLM训练进行物理-逻辑拓扑协同设计。
这些工作确立了晶圆级交换机作为独立研究方向的价值，
也揭示了该领域的核心困难：设计空间极大，且体系结构决策与多类物理约束深度耦合。
